\documentclass{beamer}
\usetheme{Warsaw}

\usepackage[T1]{fontenc}
\usepackage[utf8]{inputenc}
\usepackage{amsmath}
\usepackage{amssymb}
\usepackage{array}
\usepackage{hyperref}
\hypersetup{
    colorlinks=true,
    urlcolor=blue
}

\setbeamertemplate{footline}[frame number]

\begin{document}

\title[Stochastik]
{Wahrscheinlichkeitstheorie und Statistik (Stochastik)}
\author[Dr. Johannes Riesterer]
{Dr. rer. nat. Johannes Riesterer}
\date{}
\subject{Stochastik}

\frame{\titlepage}

\begin{frame}{Das Experiment}
\begin{block}{Ausgangssituation}
Wir werfen eine faire Münze \(n\)-mal.
Jeder einzelne Wurf hat zwei mögliche Ausgänge:
\[
K=\text{Kopf}, \qquad Z=\text{Zahl}.
\]
\end{block}

\begin{block}{Frage}
Wie kann man dieses Experiment mathematisch so beschreiben, dass man
später Wahrscheinlichkeiten und beobachtbare Größen systematisch untersuchen kann?
\end{block}
\end{frame}

\begin{frame}{Alle möglichen Abläufe}
\begin{block}{Vollständige Beschreibung}
Ein vollständiger Ablauf des Experiments ist eine Folge der Länge \(n\), zum Beispiel
\[
(K,Z,K,K,Z,\dots).
\]
\end{block}

\begin{block}{Menge aller Möglichkeiten}
Die Menge aller möglichen vollständigen Abläufe ist
\[
\Omega=\{K,Z\}^n.
\]
\end{block}

\begin{block}{Anzahl}
Da jeder der \(n\) Würfe genau zwei Möglichkeiten hat, gibt es insgesamt
\[
|\Omega|=2^n
\]
mögliche vollständige Abläufe.
\end{block}
\end{frame}

\begin{frame}{Elementarereignisse}
\begin{block}{Definition}
Ein \emph{Elementarereignis} ist genau ein einzelner vollständiger Ablauf des Experiments,
also ein einzelnes Element
\[
\omega\in\Omega.
\]
\end{block}

\begin{block}{Beispiel}
Für \(n=5\) ist
\[
\omega=(K,Z,K,K,Z)
\]
ein Elementarereignis.
\end{block}

\begin{block}{Bedeutung}
Ein Elementarereignis beschreibt den gesamten konkreten Ausgang des Experiments.
\end{block}
\end{frame}

\begin{frame}{Beispiel: \(n=3\)}
\begin{block}{Ergebnisraum}
Für \(n=3\) ist
\[
\Omega=\{K,Z\}^3.
\]
Die Elementarereignisse sind:
\[
(K,K,K),\ (K,K,Z),\ (K,Z,K),\ (K,Z,Z),
\]
\[
(Z,K,K),\ (Z,K,Z),\ (Z,Z,K),\ (Z,Z,Z).
\]
\end{block}
\end{frame}

\begin{frame}{Welche Aussagen über das Experiment interessieren uns?}
\begin{block}{Definition}
Ein \emph{Ereignis} ist eine Menge von Elementarereignissen, also eine Teilmenge von \(\Omega\):
\[
A\subseteq\Omega.
\]
\end{block}

\begin{block}{Bedeutung}
Ein Ereignis tritt ein, wenn das tatsächlich eingetretene Elementarereignis \(\omega\)
in der Menge \(A\) liegt.
\end{block}
\end{frame}

\begin{frame}{Beispiele für Ereignisse}
\begin{block}{Erster Wurf ist Kopf}
\[
A=\{\omega\in\Omega : \text{im ersten Wurf fällt Kopf}\}.
\]
\end{block}

\begin{block}{Genau zwei Köpfe}
\[
B=\{\omega\in\Omega : \text{genau zwei Köpfe auftreten}\}.
\]
\end{block}

\begin{block}{Mindestens ein Kopf}
\[
C=\{\omega\in\Omega : \text{mindestens einmal Kopf fällt}\}.
\]
\end{block}
\end{frame}

\begin{frame}{Elementarereignis oder Ereignis?}
\begin{block}{Elementarereignis}
\[
(K,Z,K)
\]
ist ein einzelner vollständiger Ablauf, also ein Elementarereignis.
\end{block}

\begin{block}{Ereignis}
\[
\{(K,K,Z),(K,Z,K),(Z,K,K)\}
\]
ist eine Menge mehrerer vollständiger Abläufe, also ein Ereignis.
\end{block}

\begin{block}{Merksatz}
Elementarereignisse sind die einzelnen Punkte von \(\Omega\), Ereignisse sind Mengen solcher Punkte.
\end{block}
\end{frame}


\begin{frame}{Beispiele zur Sprache der Ereignisse}
\begin{block}{Beispiel}
Sei \(A=\{\text{im ersten Wurf fällt Kopf}\}\) und
\(B=\{\text{im zweiten Wurf fällt Kopf}\}\).
\end{block}

\begin{block}{Interpretation}
\begin{itemize}
    \item \(A\cup B\): \text{„im ersten oder im zweiten Wurf fällt Kopf“.}
    \item \(A\cap B\): \text{„im ersten und im zweiten Wurf fällt Kopf“.}
    \item \(A^{\mathsf c}\): \text{„im ersten Wurf fällt nicht Kopf“, also Zahl.}
\end{itemize}
\end{block}
\end{frame}

\begin{frame}{Welche Mengen von Ereignissen wollen wir zulassen?}
\begin{block}{Idee}
Wir möchten mit einer Sammlung von Ereignissen arbeiten, die unter natürlichen Operationen stabil bleibt.
\end{block}

\begin{block}{Wichtig}
Wenn ein Ereignis betrachtet werden darf, dann sollen auch verwandte Ereignisse wieder betrachtet werden dürfen.
\end{block}
\end{frame}

\begin{frame}{Das sichere Ereignis soll dabei sein}
\begin{block}{Sicheres Ereignis}
Das Ereignis
\[
\Omega
\]
bedeutet: irgendein möglicher Ablauf tritt ein.
\end{block}

\begin{block}{Warum wichtig?}
Dieses Ereignis soll immer zugelassen sein, denn es tritt sicher ein.
\end{block}
\end{frame}

\begin{frame}{Das Gegenteil eines Ereignisses soll wieder erlaubt sein}
\begin{block}{Komplement}
Zu einem Ereignis \(A\subseteq\Omega\) gehört das Gegenereignis
\[
A^{\mathsf c}=\Omega\setminus A.
\]
\end{block}

\begin{block}{Bedeutung}
Wenn \(A\) bedeutet:
\[
\text{„im ersten Wurf fällt Kopf“},
\]
dann bedeutet \(A^{\mathsf c}\):
\[
\text{„im ersten Wurf fällt Zahl“.}
\]
\end{block}

\begin{block}{Wunsch}
Ist \(A\) ein erlaubtes Ereignis, dann soll auch \(A^{\mathsf c}\) ein erlaubtes Ereignis sein.
\end{block}
\end{frame}

\begin{frame}{Auch Vereinigungen von Ereignissen sollen wieder erlaubt sein}
\begin{block}{Vereinigung}
Sind \(A\) und \(B\) Ereignisse, dann bedeutet \(A\cup B\):
\[
\text{„\(A\) tritt ein oder \(B\) tritt ein“.}
\]
\end{block}

\begin{block}{Beispiel}
\(A=\{\text{erster Wurf ist Kopf}\}\),
\(B=\{\text{zweiter Wurf ist Kopf}\}\). Dann ist \(A\cup B\):
\[
\text{„im ersten oder im zweiten Wurf fällt Kopf“.}
\]
\end{block}
\end{frame}

\begin{frame}{Auch unendlich viele Vereinigungen sollen erlaubt sein}
\begin{block}{Idee}
In der Theorie möchte man nicht nur zwei oder drei Ereignisse zusammenfassen,
sondern auch Folgen von Ereignissen
\[
A_1,A_2,A_3,\dots
\]
gleichzeitig betrachten können.
\end{block}

\begin{block}{Daher}
Auch
\[
\bigcup_{n\in\mathbb N} A_n
\]
soll wieder ein erlaubtes Ereignis sein.
\end{block}
\end{frame}

\begin{frame}{Der Schnitt folgt aus den anderen Axiomen}
\begin{block}{Behauptung}
Wenn \(\mathcal A\) unter Komplement und Vereinigung stabil ist, dann ist auch der Schnitt automatisch enthalten.
\end{block}

\begin{block}{Beweis mit De Morgan}
Für beliebige \(A, B \in \mathcal A\) gilt:
\[
A \cap B = (A^{\mathsf c} \cup B^{\mathsf c})^{\mathsf c}
\]
Da \(\mathcal A\) unter Komplement stabil ist: \(A^{\mathsf c}, B^{\mathsf c} \in \mathcal A\).
Da \(\mathcal A\) unter Vereinigung stabil ist: \(A^{\mathsf c} \cup B^{\mathsf c} \in \mathcal A\).
Da \(\mathcal A\) wieder unter Komplement stabil ist: \((A^{\mathsf c} \cup B^{\mathsf c})^{\mathsf c} \in \mathcal A\).
Also: \(A \cap B \in \mathcal A\). \(\square\)
\end{block}
\end{frame}



\begin{frame}{Im Münzwurf-Beispiel ist alles besonders einfach}
\begin{block}{Einfachste Wahl}
Da \(\Omega\) endlich ist, betrachten wir einfach alle Teilmengen von \(\Omega\):
\[
\mathcal P(\Omega).
\]
\end{block}

\begin{block}{Vorteil}
Dann ist automatisch:
\begin{itemize}
    \item das sichere Ereignis dabei,
    \item das Gegenereignis eines Ereignisses wieder erlaubt,
    \item die Vereinigung von Ereignissen wieder erlaubt.
\end{itemize}
\end{block}
\end{frame}


\begin{frame}{Wie viele Möglichkeiten gibt es überhaupt?}
\begin{block}{Grundfrage}
Bevor wir Wahrscheinlichkeiten berechnen, müssen wir oft zählen, wie viele Möglichkeiten es gibt.
\end{block}

\begin{block}{Wichtig}
Dabei muss man unterscheiden:
\begin{itemize}
    \item spielt die Reihenfolge eine Rolle?
    \item darf sich etwas wiederholen?
\end{itemize}
\end{block}
\end{frame}

\begin{frame}{Zählen mit Reihenfolge}
\begin{block}{Mit Wiederholung}
Werden \(k\) Plätze nacheinander aus \(n\) Möglichkeiten besetzt und dürfen sich Einträge wiederholen, dann gibt es
\[
n^k
\]
Möglichkeiten.
\end{block}

\begin{block}{Ohne Wiederholung}
Werden \(k\) Plätze nacheinander aus \(n\) Möglichkeiten besetzt und darf nichts doppelt vorkommen, dann gibt es
\[
n(n-1)\cdots(n-k+1)=\frac{n!}{(n-k)!}
\]
Möglichkeiten.
\end{block}
\end{frame}

\begin{frame}{Zählen ohne Reihenfolge}
\begin{block}{Ohne Wiederholung}
Wählt man \(k\) Elemente aus \(n\) aus und die Reihenfolge spielt keine Rolle, dann gibt es
\[
\binom{n}{k}=\frac{n!}{k!(n-k)!}
\]
Möglichkeiten.
\end{block}

\begin{block}{Mit Wiederholung}
Wählt man \(k\) Elemente aus \(n\) aus, die Reihenfolge spielt keine Rolle und Wiederholung ist erlaubt, dann gibt es
\[
\binom{n+k-1}{k}
\]
Möglichkeiten.
\end{block}
\end{frame}

\begin{frame}{Warum brauchen wir diese Formeln hier?}
\begin{block}{Für den Münzwurf}
Ein vollständiger Ablauf von \(n\) Würfen ist eine Folge der Länge \(n\) mit zwei Möglichkeiten pro Platz.
\end{block}

\begin{block}{Daher}
Die Anzahl aller vollständigen Abläufe ist
\[
2^n.
\]
Das ist genau der Fall:
\[
\text{Reihenfolge wichtig, Wiederholung erlaubt.}
\]
\end{block}
\end{frame}

\begin{frame}{Wie viele Abläufe mit genau \(k\) Köpfen gibt es?}
\begin{block}{Idee}
Wir müssen nur entscheiden, an welchen \(k\) Stellen Kopf steht.
\end{block}

\begin{block}{Daher}
Die Anzahl der Abläufe mit genau \(k\) Köpfen ist
\[
\binom{n}{k}.
\]
\end{block}

\begin{block}{Warum?}
Wir wählen \(k\) Plätze aus \(n\) aus, an denen Kopf steht:
\[
\text{Reihenfolge unwichtig, keine Wiederholung.}
\]
\end{block}
\end{frame}

\begin{frame}{Wie berechnet man Wahrscheinlichkeiten?}
\begin{block}{Laplace-Idee}
Da die Münze fair ist, sind alle Elementarereignisse \(\omega\in\Omega\) gleich wahrscheinlich.
\end{block}

\begin{block}{Daher}
Für ein Ereignis \(A\subseteq\Omega\) ist die Wahrscheinlichkeit
\[
\mathbb P(A)=\frac{|A|}{|\Omega|}.
\]
\end{block}

\begin{block}{Hier}
Weil \(|\Omega|=2^n\), gilt also
\[
\mathbb P(A)=\frac{|A|}{2^n}.
\]
\end{block}
\end{frame}

\begin{frame}{Beispiel einer Wahrscheinlichkeit}
\begin{block}{Frage}
Wie groß ist die Wahrscheinlichkeit, dass im ersten Wurf Kopf fällt?
\end{block}

\begin{block}{Antwort}
Genau die Hälfte aller Elementarereignisse beginnt mit \(K\), also
\[
\mathbb P(\text{erster Wurf ist Kopf})=\frac{2^{n-1}}{2^n}=\frac12.
\]
\end{block}

\begin{block}{Anschaulich}
Der erste Wurf ist fair, unabhängig davon, was später passiert.
\end{block}
\end{frame}

\begin{frame}{Welche Eigenschaften soll die Wahrscheinlichkeitszuordnung haben?}
\begin{block}{Erwartung}
Eine Wahrscheinlichkeitszuordnung soll jedem erlaubten Ereignis eine Zahl zwischen \(0\) und \(1\) zuordnen.
\end{block}

\begin{block}{Wichtige Eigenschaften}
\begin{itemize}
    \item Das sichere Ereignis soll Wahrscheinlichkeit \(1\) haben.
    \item Unmögliche Ereignisse sollen Wahrscheinlichkeit \(0\) haben.
    \item Wenn zwei Ereignisse sich ausschließen, addieren sich ihre Wahrscheinlichkeiten.
\end{itemize}
\end{block}
\end{frame}

\begin{frame}{Eine Zahl zu jedem Ablauf}
\begin{block}{Neue Idee}
Oft möchte man aus einem vollständigen Ablauf nur eine bestimmte Zahl herauslesen.
\end{block}

\begin{block}{Beispiel}
Man kann jedem Ablauf die Zahl
\[
0 \text{ oder } 1
\]
zuordnen, je nachdem, ob im ersten Wurf Zahl oder Kopf gefallen ist.
\end{block}

\begin{block}{Allgemein}
Eine solche Zuordnung nimmt einen vollständigen Ablauf \(\omega\)
und liefert daraus eine Zahl.
\end{block}
\end{frame}

\begin{frame}{Die \(i\)-te Kopf-Anzeige}
\begin{block}{Definition durch Anschauung}
Für den \(i\)-ten Wurf ordnen wir jedem Ablauf \(\omega\) die Zahl
\[
1
\]
zu, wenn im \(i\)-ten Wurf Kopf gefallen ist, und die Zahl
\[
0
\]
wenn dort Zahl gefallen ist.
\end{block}

\begin{block}{Schreibweise}
Diese Zuordnung nennen wir \(X_i\):
\[
X_i(\omega)=
\begin{cases}
1,& \text{falls im \(i\)-ten Wurf Kopf fällt},\\
0,& \text{falls im \(i\)-ten Wurf Zahl fällt}.
\end{cases}
\]
\end{block}
\end{frame}

\begin{frame}{Beispiel für \(X_i\)}
\begin{block}{Beispiel}
Sei \(n=4\) und
\[
\omega=(K,Z,K,Z).
\]
Dann gilt:
\[
X_1(\omega)=1,\qquad X_2(\omega)=0,\qquad X_3(\omega)=1,\qquad X_4(\omega)=0.
\]
\end{block}

\begin{block}{Bedeutung}
\(X_i\) schaut nur auf den \(i\)-ten Wurf und ignoriert den Rest.
\end{block}
\end{frame}

\begin{frame}{Ein vollständiges Beispiel für \(n=2\)}
\[
\begin{array}{c|c|c|c}
\omega & X_1(\omega) & X_2(\omega) & S_2(\omega)=X_1(\omega)+X_2(\omega)\\
\hline
(K,K) & 1 & 1 & 2\\
(K,Z) & 1 & 0 & 1\\
(Z,K) & 0 & 1 & 1\\
(Z,Z) & 0 & 0 & 0
\end{array}
\]
\end{frame}

\begin{frame}{Was sieht man an der Tabelle?}
\begin{block}{Einzelne Anzeigen}
\(X_1\) liest nur den ersten Wurf, \(X_2\) nur den zweiten Wurf.
\end{block}

\begin{block}{Summe}
Die Größe
\[
S_2=X_1+X_2
\]
zählt die Anzahl der Köpfe.
\end{block}

\begin{block}{Verteilung}
Aus der Tabelle kann man direkt ablesen:
\[
\mathbb P(X_1=1)=\frac12,\qquad \mathbb P(X_2=1)=\frac12,
\]
\[
\mathbb P(S_2=0)=\frac14,\qquad
\mathbb P(S_2=1)=\frac12,\qquad
\mathbb P(S_2=2)=\frac14.
\]
\end{block}
\end{frame}

\begin{frame}{Aus einzelnen Anzeigen wird eine neue Größe}
\begin{block}{Idee}
Aus den einzelnen Anzeigen \(X_1,\dots,X_n\) kann man eine neue Größe bilden:
\[
S_n=X_1+\dots+X_n.
\]
\end{block}

\begin{block}{Interpretation}
Für jeden Ablauf \(\omega\) ist \(S_n(\omega)\) genau die Gesamtzahl der Köpfe.
\end{block}

\begin{block}{Mögliche Werte}
\[
0,1,2,\dots,n.
\]
\end{block}
\end{frame}

\begin{frame}{Beispiel für die Anzahl der Köpfe}
\begin{block}{Beispiel}
Für
\[
\omega=(K,Z,K,K,Z)
\]
gilt:
\[
X_1(\omega)=1,\quad X_2(\omega)=0,\quad X_3(\omega)=1,\quad X_4(\omega)=1,\quad X_5(\omega)=0.
\]
Also
\[
S_5(\omega)=3.
\]
\end{block}

\begin{block}{Anschaulich}
Die Summe zählt einfach, wie oft Kopf aufgetreten ist.
\end{block}
\end{frame}

\begin{frame}{Wie oft tritt ein bestimmter Wert auf?}
\begin{block}{Frage}
Wie oft nimmt die Zuordnung \(X_i\) den Wert \(1\) an?
\end{block}

\begin{block}{Antwort}
Genau in der Hälfte aller Elementarereignisse.
Also gilt:
\[
\mathbb P(X_i=1)=\frac12,\qquad \mathbb P(X_i=0)=\frac12.
\]
\end{block}

\begin{block}{Bedeutung}
Wir betrachten also nicht mehr nur Abläufe selbst, sondern die Verteilung
der durch \(X_i\) erzeugten Zahlen.
\end{block}
\end{frame}

\begin{frame}{Wie ist es bei der Kopfzahl \(S_n\)?}
\begin{block}{Frage}
Wie oft tritt der Wert \(k\) auf?
\end{block}

\begin{block}{Antwort}
Genau dann, wenn in dem Ablauf genau \(k\) Köpfe vorkommen.
Davon gibt es
\[
\binom{n}{k}
\]
verschiedene Elementarereignisse.
\end{block}

\begin{block}{Folgerung}
Also gilt:
\[
\mathbb P(S_n=k)=\frac{\binom{n}{k}}{2^n}.
\]
\end{block}
\end{frame}

\begin{frame}{Zwei solche Zuordnungen gleichzeitig}
\begin{block}{Neue Idee}
Man kann zwei der Zuordnungen gleichzeitig betrachten, zum Beispiel \(X_i\) und \(X_j\).
\end{block}

\begin{block}{Dann interessiert}
Nicht nur:
\[
X_i(\omega),
\]
sondern das Paar
\[
(X_i(\omega),X_j(\omega)).
\]
\end{block}

\begin{block}{Mögliche Werte}
\[
(0,0),\ (0,1),\ (1,0),\ (1,1).
\]
\end{block}
\end{frame}

\begin{frame}{Beispiel für zwei Würfe gleichzeitig}
\begin{block}{Beispiel}
Sei \(i=1\), \(j=3\) und
\[
\omega=(K,Z,Z,K).
\]
Dann gilt
\[
X_1(\omega)=1,\qquad X_3(\omega)=0,
\]
also
\[
(X_1(\omega),X_3(\omega))=(1,0).
\]
\end{block}

\begin{block}{Bedeutung}
Wir betrachten also zwei Beobachtungsgrößen gleichzeitig.
\end{block}
\end{frame}

\begin{frame}{Kombinationen von Werten}
\begin{block}{Welche Paare sind möglich?}
Für zwei Würfe gibt es genau vier mögliche Paare:
\[
(0,0),\ (0,1),\ (1,0),\ (1,1).
\]
\end{block}

\begin{block}{Bei fairer Münze}
Jedes dieser vier Paare tritt gleich häufig auf.
Also gilt:
\[
\mathbb P(X_i=0,X_j=0)=\frac14,
\]
\[
\mathbb P(X_i=0,X_j=1)=\frac14,
\]
\[
\mathbb P(X_i=1,X_j=0)=\frac14,
\]
\[
\mathbb P(X_i=1,X_j=1)=\frac14.
\]
\end{block}
\end{frame}

\begin{frame}{Getrennte Bedingungen an zwei Komponenten}
\begin{block}{Beispiel}
Man kann sagen:
\[
\text{„Die erste Größe soll \(1\) sein, die zweite darf \(0\) oder \(1\) sein.“}
\]
\end{block}

\begin{block}{Dann erhält man die Menge}
\[
\{(1,0),(1,1)\}.
\]
\end{block}

\begin{block}{Idee}
Solche Mengen beschreiben Bedingungen an die erste und zweite Komponente getrennt.
Später werden genau solche Mengen besonders wichtig.
\end{block}
\end{frame}

\begin{frame}{Unabhängigkeit anschaulich}
\begin{block}{Frage}
Beeinflusst das Ergebnis des \(i\)-ten Wurfs das Ergebnis des \(j\)-ten Wurfs?
\end{block}

\begin{block}{Beim fairen Münzwurf}
Nein. Ob im \(i\)-ten Wurf Kopf oder Zahl fällt, ändert nichts an den Möglichkeiten
des \(j\)-ten Wurfs.
\end{block}

\begin{block}{Daher}
Die Wahrscheinlichkeit für ein gemeinsames Paar zerfällt in das Produkt:
\[
\frac14=\frac12\cdot\frac12.
\]
\end{block}
\end{frame}

\begin{frame}{Ein Beispiel für Unabhängigkeit}
\begin{block}{Beispiel}
Die Wahrscheinlichkeit dafür, dass sowohl im \(i\)-ten als auch im \(j\)-ten Wurf Kopf fällt, ist
\[
\mathbb P(X_i=1,\ X_j=1)=\frac14.
\]
\end{block}

\begin{block}{Einzelwahrscheinlichkeiten}
Gleichzeitig gilt
\[
\mathbb P(X_i=1)=\frac12,\qquad \mathbb P(X_j=1)=\frac12.
\]
\end{block}

\begin{block}{Vergleich}
Also
\[
\mathbb P(X_i=1,\ X_j=1)=\mathbb P(X_i=1)\cdot\mathbb P(X_j=1).
\]
\end{block}
\end{frame}

\begin{frame}{Was haben wir bis jetzt aufgebaut?}
\begin{block}{Vom Experiment aus}
Aus dem \(n\)-fachen Münzwurf haben wir bisher ganz konkret gewonnen:
\begin{itemize}
    \item eine Menge aller Elementarereignisse,
    \item Ereignisse als Mengen von Elementarereignissen,
    \item Regeln dafür, welche Ereignisse zusammen wieder betrachtet werden dürfen,
    \item eine Regel zur Berechnung von Wahrscheinlichkeiten,
    \item Zuordnungen, die aus Elementarereignissen Zahlen machen,
    \item die Untersuchung, wie oft bestimmte Werte auftreten,
    \item die gemeinsame Betrachtung mehrerer solcher Zuordnungen,
    \item die Idee der Unabhängigkeit.
\end{itemize}
\end{block}
\end{frame}



\end{document}